% Quelle: 6. Übung CT Lösungen.pdf
% Art: Übung mit Lösungen
% Themen: irreduzible Polynome über F2, Zerlegung von X^n-1, Kreisteilungsklassen, binäre zyklische Codes, Generatorpolynom, zyklische Generatormatrix, Hammingcode H2(3), Paritätscode, ternärer Hammingcode H3(2), ternäre zyklische Codes, Standardarray-Decodierung, Syndrom, Klassenführer, [10,8,3]_11-Code, Inverse modulo 11

\section{Übung 6 (mit Lösungen)}
% Seite 1 der Quelle: Aufgabenblatt; Seiten 2--4: Lösungen.

\subsection{Binäre zyklische Codes}

\begin{aufgabe}[6\_1]
\begin{enumerate}[label=\alph*)]
  \item Bestimmen Sie sämtliche irreduzible Polynome 3.\ Grades aus $F_2[X]$.
  \item Zerlegen Sie $X^7 - 1 \in F_2[X]$ in irreduzible Polynome.
\end{enumerate}
\end{aufgabe}

\begin{loesung}
% KORRIGIERT: Begründung fehlte (Absolutglied 1; Grad 3: irreduzibel genau dann, wenn keine Nullstelle)
\textbf{a)} Ein Polynom vom Grad 3 über $F_2$ ist genau dann irreduzibel, wenn es keine Nullstelle in $F_2$ hat.
Wegen $X = 0$ muss das Absolutglied $1$ sein. Ein solches Polynom hat also die Form $X^3 + aX^2 + bX + 1$, $a, b \in F_2$.\\
Für $a = b = 0$ oder $a = b = 1$ ist $X = 1$ Nullstelle (denn $1 + a + b + 1 \equiv a + b \equiv 0 \pmod{2}$), also $a \neq b$.\\
$\Rightarrow\ X^3 + X + 1$ oder $X^3 + X^2 + 1$

\medskip
% KORRIGIERT: Original ohne Malpunkte, ohne Modul und ohne Angabe von n, p: "{s, sp, sp^2, ..., sp^{d-1}}"
\textbf{b)} Kreisteilungsklassen modulo $n = 7$ zu $p = 2$: $\{s, s \cdot p, s \cdot p^2, \dots, s \cdot p^{d-1}\} \pmod{7}$, wobei $d$ die kleinste Zahl mit $s \cdot p^d \equiv s \pmod{7}$ ist\\
$\{0\}$, $\{1, 2, -3\}$, $\{-1, -2, 3\}$ \quad (wegen $4 \equiv -3$, $6 \equiv -1$, $5 \equiv -2 \pmod{7}$)\\
$X^7 - 1 = (X + 1)(X^3 + X + 1)(X^3 + X^2 + 1)$
\end{loesung}

\begin{aufgabe}[6\_2]
Bestimmen Sie vom binären Hamming Code $H_2(3)$
\begin{itemize}
  \item ein Generatorpolynom $g(X)$,
  \item eine zyklische Generatormatrix $G$.
\end{itemize}
Hinweis: Es gibt zwei verschiedene, aber äquivalente solche Codes.
\end{aufgabe}

\begin{loesung}
Der binäre Hamming Code $H_2(3)$ hat die Parameter $[7, 4, 3]_2$,\\
d.h.\ Wortlänge $n = 7$, $k = \dim C = 4$, $d_C = 3$.

Generatorpolynom $g(X) \mid X^7 - 1$, \quad $\operatorname{grad} g(X) = p = n - k = 3$.\\
Zwei verschiedene Generatorpolynome führen auf äquivalente zyklische $[7, 4]_2$ Hammingcodes.

\medskip
1.\ Fall: \quad $g(X) = 1 + X + X^3 \quad\Rightarrow\quad
G = \begin{pmatrix}
1 & 1 & 0 & 1 & 0 & 0 & 0\\
0 & 1 & 1 & 0 & 1 & 0 & 0\\
0 & 0 & 1 & 1 & 0 & 1 & 0\\
0 & 0 & 0 & 1 & 1 & 0 & 1
\end{pmatrix}$

\medskip
2.\ Fall: \quad $g(X) = 1 + X^2 + X^3 \quad\Rightarrow\quad
G = \begin{pmatrix}
1 & 0 & 1 & 1 & 0 & 0 & 0\\
0 & 1 & 0 & 1 & 1 & 0 & 0\\
0 & 0 & 1 & 0 & 1 & 1 & 0\\
% KORRIGIERT: letzte Zeile lautete "0 0 0 1 1 0 1"; die Verschiebung X^3 g(X) = X^3 + X^5 + X^6 ergibt "0 0 0 1 0 1 1"
0 & 0 & 0 & 1 & 0 & 1 & 1
\end{pmatrix}$
\end{loesung}

\begin{aufgabe}[6\_3]
Bestimmen Sie von dem binären $[7, 6, 2]_2$ Paritätscode
\begin{itemize}
  \item das Generatorpolynom $g(X)$,
  \item eine zyklische Generatormatrix $G$.
\end{itemize}
\end{aufgabe}

\begin{loesung}
% KORRIGIERT: "[7, 6]" ohne Alphabetgröße
Der $[7, 6]_2$ Paritätscode hat Wortlänge $n = 7$ und $k = \dim C = 6$.

Generatorpolynom $g(X) = X + 1 \mid X^7 - 1$. \quad $\operatorname{grad} g(X) = p = n - k = 1$
\[
G = \begin{pmatrix}
1 & 1 & 0 & 0 & 0 & 0 & 0\\
0 & 1 & 1 & 0 & 0 & 0 & 0\\
0 & 0 & 1 & 1 & 0 & 0 & 0\\
0 & 0 & 0 & 1 & 1 & 0 & 0\\
0 & 0 & 0 & 0 & 1 & 1 & 0\\
0 & 0 & 0 & 0 & 0 & 1 & 1
\end{pmatrix}
\]
\end{loesung}

\subsection{Ternäre Codes mit den Parametern $[n, k, d]_3 = [4, 2, d]_3$}

\begin{aufgabe}[6\_4]
Erinnern Sie sich an: die Parameter $[n, k, d]$, eine Generatormatrix $G$ und eine Kontrollmatrix $H$ des ternären Hammingcodes $H_3(2)$.
\end{aufgabe}

\begin{loesung}
Ternärer Hammingcode $H_3(2)$\\
Parameter $[4, 2, 3]_3$
% KORRIGIERT: "=" statt "≡ (mod 3)" (Einträge von -A^T modulo 3 reduziert)
\[
H = (\, E \mid A \,) = \begin{pmatrix} 1 & 0 & 1 & 1\\ 0 & 1 & 1 & 2 \end{pmatrix}, \qquad
G = (-A^T \mid E\,) \equiv \begin{pmatrix} 2 & 2 & 1 & 0\\ 2 & 1 & 0 & 1 \end{pmatrix} \pmod{3}
\]
\end{loesung}

\begin{aufgabe}[6\_5]
Zerlegen Sie $X^4 - 1 \in F_3[X]$ in irreduzible Polynome.
\end{aufgabe}

\begin{loesung}
% KORRIGIERT: Original ohne Malpunkte und ohne Modul: "{s, sp, sp^2, ..., sp^{d-1}}"
Kreisteilungsklassen modulo $n = 4$ zu $p = 3$: $\{s, s \cdot p, s \cdot p^2, \dots, s \cdot p^{d-1}\} \pmod{4}$\\
$\{0\}$, $\{1, -1\}$, $\{2\}$ \quad (wegen $3 \equiv -1 \pmod{4}$)
\begin{align*}
X^4 - 1 &= (X - 1)(X^3 + X^2 + X + 1) = (X - 1)\bigl(X^2(X + 1) + (X + 1)\bigr)\\
        &= (X - 1)(X + 1)(X^2 + 1)
\end{align*}
% KORRIGIERT: Begründung der Irreduzibilität von X^2 + 1 fehlte
$X^2 + 1$ ist irreduzibel über $F_3$, denn es hat Grad 2 und keine Nullstelle: $0^2 + 1 = 1$, $1^2 + 1 = 2$, $2^2 + 1 = 5 \equiv 2 \pmod{3}$.
\end{loesung}

\begin{aufgabe}[6\_6]
Es gibt 2 zyklische ternäre Codes mit den Parametern $n = 4$, $k = 2$, welche?
Bestimmen Sie dazu $d$, $g(X)$, $G$ dieser Codes.
% Im Lösungsteil lautet die Aufgabenstellung: "Bestimmen Sie dazu d, g(X), h(X), G, H dieser Codes."
\end{aufgabe}

\begin{loesung}
(Aufgabenstellung im Lösungsteil: Bestimmen Sie dazu $d$, $g(X)$, $h(X)$, $G$, $H$ dieser Codes.)
\begin{enumerate}
  \item $g(X) = X^2 + 1$, \quad $G = \begin{pmatrix} 1 & 0 & 1 & 0\\ 0 & 1 & 0 & 1 \end{pmatrix}$, \quad $d = 2$, $e = 0$
  \item $g(X) = X^2 - 1$, \quad $G = \begin{pmatrix} -1 & 0 & 1 & 0\\ 0 & -1 & 0 & 1 \end{pmatrix}$, \quad $d = 2$, $e = 0$
\end{enumerate}
\end{loesung}

\begin{aufgabe}[6\_7]
Gibt es folglich einen zu $H_3(2)$ isomorphen zyklischen Code?
\end{aufgabe}

\begin{loesung}
Nein, denn beide zyklischen Codes haben nur $d = 2$, % KORRIGIERT: "also d = 3" -- aus "korrigiert einen Fehler" folgt nur d >= 3; d = 3 ist der Parameter aus 6_4
während $H_3(2)$ einen Fehler korrigieren kann, also $d \geq 3$ (genauer $d = 3$, siehe Aufgabe 6\_4). Äquivalente Codes haben dieselbe Minimaldistanz.
\end{loesung}

\subsection{Standardarray-Decodierung}

\begin{aufgabe}[6\_8]
Sei $C = H_3(2)$ der ternäre Hammingcode aus Aufgabe 6\_4. Bestimmen Sie
\begin{itemize}
  \item die Klassenführer,
  \item die Syndrome der Nebenklassen und
  \item dekodieren Sie $r_1 = 2\,2\,0\,0$, $r_2 = 0\,1\,2\,1$.
\end{itemize}
\end{aufgabe}

\begin{loesung}
Klassenführer, Syndrome der Nebenklassen, $\mathbf{r}_1 = 2\,2\,0\,0$, $\mathbf{r}_2 = 0\,1\,2\,1$
\[
H = (\, E \mid A \,) = \begin{pmatrix} 1 & 0 & 1 & 1\\ 0 & 1 & 1 & 2 \end{pmatrix}
\]
\begin{center}
\begin{tabular}{ll}
\toprule
Klassenführer & Syndrom\\
\midrule
0 0 0 0 & 0 0\\
1 0 0 0 & 1 0\\
0 1 0 0 & 0 1\\
0 0 1 0 & 1 1\\
0 0 0 1 & 1 2\\
2 0 0 0 & 2 0\\
0 2 0 0 & 0 2\\
0 0 2 0 & 2 2\\
0 0 0 2 & 2 1\\
\bottomrule
\end{tabular}
\end{center}

$\operatorname{Syn}(\mathbf{r}) = H\mathbf{r}^T$

\medskip
$\mathbf{r}_1 = 2\,2\,0\,0$
\[
% KORRIGIERT: fehlende Malpunkte, fehlendes (mod 3) in Syndrom- und Codewortrechnung (6_8)
\operatorname{Syn} \mathbf{r}_1 = H\mathbf{r}_1^T = 2 \cdot \begin{pmatrix} 1\\ 0 \end{pmatrix} + 2 \cdot \begin{pmatrix} 0\\ 1 \end{pmatrix} = \begin{pmatrix} 2\\ 2 \end{pmatrix}
\quad\Rightarrow \text{Klassenführer } \mathbf{e}_1 = 0\,0\,2\,0
\]
d.h.\ Fehlerstelle $x = 3$, Fehlergröße $e_x = 2$\\
$\mathbf{c}_1 = \mathbf{r}_1 - \mathbf{e}_1 = 2\,2\,(-2)\,0 \equiv 2\,2\,1\,0 \pmod{3}$

\medskip
$\mathbf{r}_2 = 0\,1\,2\,1$
\[
\operatorname{Syn} \mathbf{r}_2 = H\mathbf{r}_2^T = 1 \cdot \begin{pmatrix} 0\\ 1 \end{pmatrix} + 2 \cdot \begin{pmatrix} 1\\ 1 \end{pmatrix} + 1 \cdot \begin{pmatrix} 1\\ 2 \end{pmatrix} = \begin{pmatrix} 3\\ 5 \end{pmatrix} \equiv \begin{pmatrix} 0\\ 2 \end{pmatrix} \pmod{3}
\quad\Rightarrow \text{Klassenführer } \mathbf{e}_2 = 0\,2\,0\,0
\]
d.h.\ Fehlerstelle $x = 2$, Fehlergröße $e_x = 2$\\
$\mathbf{c}_2 = \mathbf{r}_2 - \mathbf{e}_2 = 0\,(-1)\,2\,1 \equiv 0\,2\,2\,1 \pmod{3}$
\end{loesung}

\subsection{$[10, 8, 3]_{11}$ Code}

\begin{aufgabe}[6\_9]
Decodieren Sie die empfangenen Wörter
\begin{align*}
r_1 &= 1\,1\,5\,1\,1\,1\,1\,1\,0\,3\\
r_2 &= 1\,1\,5\,6\,1\,1\,1\,1\,0\,3\\
r_3 &= 1\,1\,1\,1\,1\,1\,1\,1\,0\,3
\end{align*}
\end{aufgabe}

\begin{loesung}
Inverse Modulo 11
\begin{center}
% KORRIGIERT: Tabelle ohne Hinweis, dass a * a^{-1} ≡ 1 (mod 11) gemeint ist
(Es gilt $a \cdot a^{-1} \equiv 1 \pmod{11}$; alle Inversen existieren, da $11$ prim ist und $\ggT(a, 11) = 1$ für $a \not\equiv 0$.)
\begin{tabular}{c|rrrrrrrrrr}
$a$      & 1 & 2  & 3 & 4 & 5  & $-5$ & $-4$ & $-3$ & $-2$ & $-1$\\
\midrule
$a^{-1}$ & 1 & $-5$ & 4 & 3 & $-2$ & 2 & $-3$ & $-4$ & 5 & $-1$\\
\end{tabular}
\end{center}

\begin{align*}
\mathbf{r}_1 &= 1\,1\,5\,1\,1\,1\,1\,1\,0\,3\\
\mathbf{r}_2 &= 1\,1\,5\,6\,1\,1\,1\,1\,0\,3\\
\mathbf{r}_3 &= 1\,1\,1\,1\,1\,1\,1\,1\,0\,3
\end{align*}

% KORRIGIERT: "≡" ohne Modulangabe
\[
H = \begin{pmatrix}
1 & 1 & 1 & 1 & 1 & 1 & 1 & 1 & 1 & 1\\
1 & 2 & 3 & 4 & 5 & 6 & 7 & 8 & 9 & 10
\end{pmatrix}
\equiv
\begin{pmatrix}
1 & 1 & 1 & 1 & 1 & 1 & 1 & 1 & 1 & 1\\
1 & 2 & 3 & 4 & 5 & -5 & -4 & -3 & -2 & -1
\end{pmatrix}
\pmod{11}
\]

% KORRIGIERT: Fehlergröße hieß erst "m", danach "y"; "=" statt "≡ (mod 11)"; Voraussetzungen (ein Fehler, y ≠ 0) fehlten
Bei genau einem Fehler an Stelle $x$ mit Fehlergröße $y \not\equiv 0$ gilt:
$\mathbf{e} = \mathbf{r} - \mathbf{c} = (\,0 \dots y \dots 0\,)$, \quad
$S(\mathbf{r}) = H\mathbf{r}^T = \begin{pmatrix} s_0\\ s_1 \end{pmatrix} \equiv y \cdot \begin{pmatrix} 1\\ x \end{pmatrix} \pmod{11}$, \quad
also $y \equiv s_0$ und $x \equiv s_1 \cdot s_0^{-1} \pmod{11}$ (die Inverse existiert, da $s_0 \not\equiv 0$), \quad
$\mathbf{c} = \mathbf{r} - \mathbf{e}$

\[
% KORRIGIERT: "=" statt "≡ (mod 11)" in den Syndromrechnungen (s_1(r_1) = 78 ≡ 1 ≡ 4 · 3)
S(\mathbf{r}_1) = \begin{pmatrix} 15\\ 78 \end{pmatrix} \equiv \begin{pmatrix} 4\\ 4 \cdot 3 \end{pmatrix} = 4 \cdot \begin{pmatrix} 1\\ 3 \end{pmatrix} \pmod{11}
\qquad\Rightarrow\ x = 3,\ y = 4
\]
$\Rightarrow \mathbf{c}_1 = \mathbf{r}_1 - \mathbf{e} = 1\,1\,(5 - 4)\,1\,1\,1\,1\,1\,0\,3 = 1\,1\,1\,1\,1\,1\,1\,1\,0\,3 = \mathbf{r}_3$

\[
S(\mathbf{r}_2) \equiv \begin{pmatrix} 4 + 5\\ 1 + 4 \cdot 5 \end{pmatrix} = \begin{pmatrix} 9\\ 21 \end{pmatrix} \equiv \begin{pmatrix} -2\\ -1 \end{pmatrix} \equiv (-2) \cdot \begin{pmatrix} 1\\ 6 \end{pmatrix} \pmod{11}
\qquad\Rightarrow\ y = -2,\ x \equiv (-1) \cdot (-2)^{-1} \equiv (-1) \cdot 5 \equiv 6 \pmod{11}
\]
($\mathbf{r}_2$ unterscheidet sich von $\mathbf{r}_1$ nur an Stelle $4$ um $+5$, daher $S(\mathbf{r}_2) = S(\mathbf{r}_1) + 5 \cdot (1, 4)^T$.)
$\Rightarrow \mathbf{c}_2 = \mathbf{r}_2 - \mathbf{e} = 1\,1\,5\,6\,1\,(1 + 2)\,1\,1\,0\,3 = 1\,1\,5\,6\,1\,3\,1\,1\,0\,3$

\[
S(\mathbf{r}_3) = \begin{pmatrix} 11\\ 66 \end{pmatrix} \equiv \begin{pmatrix} 0\\ 0 \end{pmatrix} \pmod{11} \qquad\Rightarrow \mathbf{c}_3 = \mathbf{r}_3
\]
\end{loesung}
